\appendix
\section{Proofs of every result used}
This appendix makes the paper standalone. Each result used in the text and not already proved where it is stated is restated and proved here, including those taken from the author's earlier papers. Where a proof rests on an outside input, the input is named and is the only thing assumed. Items marked \textsf{[Lean]} are also checked by machine (Appendix B).

\paragraph{A.1 Zero and span are free (§\ref{sec:size}).}
\emph{Statement.} For a Hermitian generator $K$, a shift $K\mapsto K+cI$ leaves every eigenvector and every eigenvalue difference unchanged, and a positive rescaling $K\mapsto cK$, $c>0$, leaves every eigenvector, the order of the eigenvalues, every spectral projector $\mathbf 1[K<0]$, and so the ground-state correlation matrix and its census, unchanged.
\emph{Proof.} $K+cI$ commutes with $K$ and has eigenvalues $\lambda_i+c$, so $\lambda_i-\lambda_j$ is unchanged. $cK$ has the same eigenvectors with eigenvalues $c\lambda_i$; since $c>0$, $c\lambda_i<0\iff\lambda_i<0$, so the filled projector $P=\mathbf 1[K<0]$ is the same operator. The regional correlation matrix $C_A=E^\dagger PE$ and every function of its spectrum, including the census, are unchanged; so are principal angles between spectral subspaces, which depend only on the subspaces. Under a shift the filled projector is $\mathbf 1[K<-c]$ for the same Fermi level, which is why the sector must be carried with the shift. \qed

\paragraph{A.2 The census coefficient $A=2c/\pi^2$ (§\ref{sec:size}, eq.~\ref{eq:census}).}
\emph{Inputs.} (i) The single-particle entanglement spectrum at one cut is equally spaced, $\varepsilon_j=(j+\tfrac12)\Delta$ on both signs, as measured on the Ising chain (script \texttt{s02}) and derived in \emph{The Angle Between the Rails}. (ii) The entanglement entropy of one cut in a gapped chain near criticality is $S_{\mathrm{cut}}=\tfrac c6\ln(\xi/a)+O(1)$ (Calabrese–Cardy).
\emph{Derivation.} A level is in the window $\nu\in[\epsilon,1-\epsilon]$ iff $|\varepsilon|\le\Lambda=\ln((1-\epsilon)/\epsilon)$ (Proposition~\ref{prop:bridge}). With spacing $\Delta$ on both signs, one cut has $N_{\mathrm{cut}}=2\Lambda/\Delta+O(1)$ levels in the window. Its entropy is $S_{\mathrm{cut}}=\sum_jH(\nu(\varepsilon_j))=\Delta^{-1}\int H\,d\varepsilon+O(\Delta)=\pi^2/(3\Delta)+O(\Delta)$ by Proposition~\ref{prop:bridge} (midpoint rule). Equating with (ii), $1/\Delta=(c/2\pi^2)\ln(\xi/a)$, so $N_{\mathrm{cut}}=(c/\pi^2)\ln(\xi/a)\,\Lambda$, and an interval with two cuts has $A=2c/\pi^2$. \qed\ The derivation is exact given (i) and (ii); its corrections are those of the two inputs, and the level-counting convention (one complex mode per level) must match the entropy convention used in (ii).

\paragraph{A.3 The bridge, and $\int H\,d\varepsilon=\pi^2/3$ (Proposition~\ref{prop:bridge}).}
Proved in the text. The last step uses $\int_0^1(-\ln\nu)/(1-\nu)\,d\nu=\sum_{n\ge0}\int_0^1\nu^n(-\ln\nu)\,d\nu=\sum_{n\ge1}n^{-2}=\pi^2/6$, and $\int_0^1H/(\nu(1-\nu))\,d\nu=2\int_0^1(-\ln\nu)/(1-\nu)\,d\nu$ by the symmetry $\nu\leftrightarrow1-\nu$ and $H/(\nu(1-\nu))=-\ln\nu/(1-\nu)-\ln(1-\nu)/\nu$. The identities $\Gamma=1-2\nu$ and $\nu(1-\nu)=(1-\Gamma^2)/4$ \textsf{[Lean: \texttt{bridge}, \texttt{census\_square}]}.

\paragraph{A.4 The offset is the log-odds of empty against full (§\ref{sec:ledger}; \emph{The Offset Belongs to the Boundary}, Thm 2).}
\emph{Statement.} For a number-conserving Gaussian fermion state with $0<C_A<I$, $\tau=\Tr[\log(I-C_A)-\log C_A]=\log\det(I-C_A)-\log\det C_A=\log P(N_A=0)/P(N_A=|A|)$.
\emph{Proof.} The reduced state of a Gaussian state is Gaussian, $\rho_A\propto\exp(-\sum_k\varepsilon_kb_k^\dagger b_k)$ in the eigenmodes $b_k$ of $C_A$, a product over modes with occupation $\nu_k=1/(1+e^{\varepsilon_k})$. The region is empty iff every mode is empty and full iff every mode is filled, so $P(0)=\prod_k(1-\nu_k)=\det(I-C_A)$ and $P(|A|)=\prod_k\nu_k=\det C_A$. Taking logs, $\tau=\sum_k\ln((1-\nu_k)/\nu_k)=\sum_k\varepsilon_k$. The example in the text: $\mathrm{diag}(0.2,0.8)$ gives $\ln4+\ln\tfrac14=0$, $\mathrm{diag}(0.2,0.2)$ gives $2\ln4=\ln16$, while both have $C-C^2=0.16I$. \qed

\paragraph{A.5 Defect split for the correlation matrix (Theorem~\ref{thm:weight}(1)).}
Block multiplication gives $(C^2)_{AA}=C_A^2+C_{AB}C_{BA}$ \textsf{[Lean: \texttt{topleft\_sq}; formal suite: \texttt{block\_defect\_identity}]}; subtracting it from $C_A$ gives the split. If $C^2=C$ the first term vanishes \textsf{[Lean: \texttt{idempotent\_cross}]}. For Hermitian $C$, $C_{BA}=C_{AB}^\dagger$, so $\Tr C_{AB}C_{AB}^\dagger=\|C_{AB}\|_F^2$, and $P_AC(1-P_A)$ is $C_{AB}$ padded with zeros. The same functional applied to a generator $H$ and a moving projector $P$ is $E_H=\|PH(1-P)\|_F^2$. \qed

\paragraph{A.6 The boundary coordinate is the modulus, and it survives opening (§\ref{sec:cross}, §\ref{sec:opened}).}
\emph{Statement.} If $C^2=C$, each eigenvalue $\lambda$ of $X_A$ is $\nu(1-\nu)$ for an eigenvalue $\nu$ of $C_A$, and $\sqrt{|1-4\lambda|}=|1-2\nu|$. This holds for the biorthogonal projector $C=\sum_{k\,\mathrm{filled}}|R_k\rangle\langle L_k|$ with $\langle L_j|R_k\rangle=\delta_{jk}$.
\emph{Proof.} By A.5, $X_A=C_A-C_A^2=f(C_A)$ with $f(x)=x-x^2$, so by the spectral mapping theorem (valid for any square matrix, via the Jordan form) its eigenvalues are $f(\nu)$. Then $1-4f(\nu)=(1-2\nu)^2$ \textsf{[Lean: \texttt{census\_square}]}, whose modulus has square root $|1-2\nu|$ for complex $\nu$ too. For the biorthogonal projector, $C^2=\sum_{j,k}|R_j\rangle\langle L_j|R_k\rangle\langle L_k|=C$. \qed

\paragraph{A.7 The Offset Speed Limit (Theorem~\ref{thm:speed}; \emph{The Offset Speed Limit}, Thm 1 and Cor.~3).}
\emph{Proof.} \emph{Step 1 (trace derivative).} $\frac{d}{dt}\Tr\log(I-C)=-\Tr[(I-C)^{-1}\dot C]$ and $\frac d{dt}\Tr\log C=\Tr[C^{-1}\dot C]$; since $(I-C)^{-1}+C^{-1}=(C(I-C))^{-1}$, $\dot\tau=\Tr[W\dot C]$. \emph{Step 2 (lift).} $\dot C=E^\dagger\dot PE$, so $\dot\tau=\Tr[EWE^\dagger\dot P]=\Tr[F\dot P]$, with $F$ Hermitian. \emph{Step 3 (only the off-diagonal part counts).} Differentiating $P^2=P$ gives $\dot PP+P\dot P=\dot P$, hence $P\dot PP=0$ and $(1-P)\dot P(1-P)=0$: $\dot P$ is block off-diagonal. So $\Tr[F\dot P]=\Tr[F_{\mathrm{off}}\dot P]$ with $F_{\mathrm{off}}=PF(1-P)+(1-P)FP$, and expanding, $[P,[P,F]]=PF+FP-2PFP=F_{\mathrm{off}}=T$. \emph{Step 4 (Cauchy–Schwarz).} In the Hilbert–Schmidt inner product, $\dot\tau^2=\langle T,\dot P\rangle^2\le\|T\|_F^2\|\dot P\|_F^2=2g\|T\|_F^2$, with equality iff $\dot P=\lambda T$, $\lambda\in\mathbb R$. \emph{Step 5 ($\|T\|_F^2=2S$).} $\|T\|_F^2=2\Tr[PF(1-P)FP]=2\Tr[W\,E^\dagger(1-P)E\,W\,E^\dagger PE]=2\Tr[W(I-C)WC]=2\Tr[W^2C(I-C)]=2S$, because $W$, $C$ and $I-C$ commute. Hence $\dot\tau^2\le4gS$. For the corollary, $[H,P]=(1-P)HP-PH(1-P)$, whose two blocks are adjoint up to sign, so $\|\dot P\|_F^2=2\|PH(1-P)\|_F^2$ and $g=E_H$. \qed

\paragraph{A.8 Peel steps add like velocities (Theorem~\ref{thm:add}; VOTKP Thm 15.1).}
Proved in the text. In the division-free form used by the Lean file, $g(z_1)+g(z_2)=g(z_1z_2)\,(1+g(z_1)g(z_2))$ for $z_1,z_2>0$, with $g(z)=(z-1)/(z+1)$ \textsf{[Lean: \texttt{rapidities\_add}]}. Since $1+g_1g_2>0$ for $|g_i|<1$, this is $g(z_1z_2)=g_1\oplus g_2$. The ladder reconstruction is the inverse $z=(1+\Gamma)/(1-\Gamma)$ applied step by step \textsf{[formal suite: \texttt{gamma\_inverse}, \texttt{gamma\_anchor\_reconstruction}]}.

\paragraph{A.9 Scale blindness (§\ref{sec:atom}; VOTKP Thm 15.4).}
\emph{Statement.} For a nonrelativistic Coulomb system $H_\lambda=\sum_i p_i^2/(2\lambda m_i)+\sum_{i<j}q_iq_j/|r_i-r_j|$, multiplying every mass by $\lambda>0$ multiplies every energy by $\lambda$, so every ratio of removal energies, and every $\Gamma$, is unchanged.
\emph{Proof.} Let $U$ be the unitary dilation $(U\psi)(r)=\lambda^{3N/2}\psi(\lambda r)$. Then $U^\dagger r_iU=r_i/\lambda$ and $U^\dagger p_iU=\lambda p_i$, so $U^\dagger H_\lambda U=\sum_i\lambda^2p_i^2/(2\lambda m_i)+\lambda\sum_{i<j}q_iq_j/|r_i-r_j|=\lambda H_1$. Unitarily equivalent operators have the same spectrum, so $\mathrm{spec}\,H_\lambda=\lambda\,\mathrm{spec}\,H_1$, and so do the thresholds of every ion, which are bottoms of spectra of subsystems. Ratios cancel $\lambda$. \qed\ Mass \emph{ratios} still enter, which is why the ${}^1$H/${}^2$H isotope shift is visible (VOTKP §15.3).

\paragraph{A.10 The last peel step at first order (§\ref{sec:atom}; VOTKP §15.10).}
\emph{Statement.} To first order in the electron repulsion, a helium-like ion of charge $Z\ge2$ has removal threshold $I_{\mathrm{He}}=Z^2/2-5Z/8$ Hartree, the hydrogen-like ion $I_{\mathrm H}=Z^2/2$, and the last step has $\Gamma=5/(8Z-5)\to5/(8Z)$.
\emph{Proof.} The hydrogen-like ground energy is $-Z^2/2$. Two non-interacting $1s$ electrons have $-Z^2$. The first-order repulsion is $J=\iint\rho(r_1)\rho(r_2)/|r_1-r_2|$ with $\rho=(Z^3/\pi)e^{-2Zr}$. By Gauss's law the potential of $\rho$ is $\phi(r)=1/r-(Z+1/r)e^{-2Zr}$, so
\[
J=4Z^3\!\int_0^\infty\! r^2e^{-2Zr}\phi(r)\,dr=4Z^3\Big[\tfrac1{4Z^2}-\tfrac{2Z}{(4Z)^3}-\tfrac1{(4Z)^2}\Big]=4Z^3\cdot\tfrac{10}{64Z^2}=\tfrac58Z
\]
(Unsöld 1927). So $E_{\mathrm{He}}=-Z^2+5Z/8$ and $I_{\mathrm{He}}=-Z^2/2-E_{\mathrm{He}}=Z^2/2-5Z/8$, positive for $Z\ge2$. With $z=I_{\mathrm H}/I_{\mathrm{He}}$, $\Gamma=(I_{\mathrm H}-I_{\mathrm{He}})/(I_{\mathrm H}+I_{\mathrm{He}})=(5Z/8)/(Z^2-5Z/8)=5/(8Z-5)$ \textsf{[Lean: \texttt{last\_step}]}. \qed\ The higher terms of the $1/Z$ series (Hylleraas; Scherr–Knight) and the Thomas–Fermi middle of the peel are cited inputs in VOTKP and are not used here.

\paragraph{A.11 Ideal Saha gas: the electron density cancels (§\ref{sec:gas}).}
\emph{Input.} The ideal Saha equation $n_{i+1}n_e/n_i=(2g_{i+1}/g_i)\lambda_e^{-3}e^{-\beta I_i}$ with constant statistical weights.
\emph{Statement.} $n_0n_2/n_1^2=(g_0g_2/g_1^2)\,e^{-\beta(I_1-I_0)}$ and $n_0/(n_2n_e^2\lambda_e^6)=(g_0/4g_2)\,e^{\beta(I_0+I_1)}$; their $\beta$-slopes return $I_1-I_0$ and $I_0+I_1$, hence $\Gamma_0=(I_1-I_0)/(I_1+I_0)$.
\emph{Proof.} Divide the $i=1$ equation by the $i=0$ one: $n_e$ and $\lambda_e$ cancel \textsf{[Lean: \texttt{saha\_ratio}]}. Multiply the two equations for the second form, which needs $n_e\lambda_e^3$ per sample. The ratio of the two slopes gives $\Gamma_0$. \qed\ This is the $Z=2$ case of Theorem~\ref{thm:inverse}.

\paragraph{A.12 Outside the rim, exactly, and the Cayley pole (§\ref{sec:nucleus}).}
For $S>s>0$ and $z=S/(-s)$: $\Gamma=(S+s)/(S-s)$ and $\Gamma-1=2s/(S-s)$; $S_2/S_1=-1\iff S_1+S_2=0$ \textsf{[Lean: \texttt{outside\_rim}, \texttt{rim\_distance}, \texttt{cayley\_pole}]}. The half-turn statement is Theorem~\ref{thm:cayley}.

\paragraph{A.13 No cone along a Hermitian family (§\ref{sec:opened}; cone paper v8, Thm 1).}
\emph{Statement.} For a two-level family $h(\theta)$ that is Hermitian along the direction of differentiation, with $E_+\ne E_-$, $Q=M_{+-}M_{-+}/(E_+-E_-)^2\ge0$, where $M_{ab}=\langle L_a|\partial_\theta h|R_b\rangle$.
\emph{Proof.} At a Hermitian point $L_a=R_a$; with $\partial_\theta h$ Hermitian, $M_{-+}=\overline{M_{+-}}$, and $E_\pm$ are real. Hence $Q=|M_{+-}|^2/(E_+-E_-)^2\ge0$. \qed\ The hypothesis on the tangent is needed: $h(\theta)=\sigma_z+i\theta\sigma_x$ is Hermitian at $\theta=0$, $\partial_\theta h=i\sigma_x$ is not, and $Q=(i)(i)/4=-\tfrac14$.

\paragraph{A.14 The Petermann factor and the exit line (Proposition~\ref{prop:petermann}).}
The formula $K=\tfrac12(1+\sum|d_j|^2/|E|^2)$ follows from $P=\tfrac12(I-h/E)$ and $\Tr\sigma_j\sigma_k=2\delta_{jk}$: $\Tr P^\dagger P=\tfrac14\Tr[(I-\bar d\cdot\sigma/\bar E)(I-d\cdot\sigma/E)]=\tfrac12(1+\bar d\cdot d/|E|^2)$, the cross terms being traceless. The final equivalence $1+u^2/p=2\iff p=u^2$ for $p>0$ \textsf{[Lean: \texttt{exit\_line}]}. The integral $\frac1{2\pi}\int_0^{2\pi}dk/(A+B\cos k)=(A^2-B^2)^{-1/2}$ for $A>|B|$ is the residue at the root of $Bz^2+2Az+B$ inside the unit circle, or the substitution $t=\tan(k/2)$ on the two half-intervals.

\paragraph{A.15 The quarter-turn and its branch (§\ref{sec:opened}, eq.~\ref{eq:quarter}).}
Let $p=E(0)E(\pi)>0$ and $\gamma=u/\sqrt p$. For real $y$, $\operatorname{artanh}(iy)=i\arctan y$, because $\tanh(i\phi)=i\tan\phi$. So $\tau_{\mathrm{end}}=2i\arctan\gamma$, $\Ima\tau_{\mathrm{end}}=\theta=2\arctan\gamma\in(0,\pi)$, and $\Gamma_{\mathrm{end}}=\tanh(i\arctan\gamma)=i\gamma$. Then $\tan^2(\theta/2)=\gamma^2=u^2/p$, so $\sec^2(\theta/2)=1+u^2/p=\langle K\rangle$, and $\langle K\rangle=2\iff\gamma=1\iff\theta=\pi/2\iff\Gamma_{\mathrm{end}}=i$. The continuation of the closed form to the opened chain is a definition; §\ref{sec:opened} tests it and finds it does not describe the census rotation away from the exit. \qed

\paragraph{A.16 The bilinear Gram determinant (Theorem~\ref{thm:signature}).}
For complex vectors $d,\alpha,\beta\in\mathbb C^3$ with the bilinear product, the Gram matrix $G$ of $(d,\alpha,\beta)$ is $M^TM$ with $M=(d\,|\,\alpha\,|\,\beta)$, so $\det G=(\det M)^2=[d\cdot(\alpha\times\beta)]^2$. Expanding $\det G$ along its first row and using $s=d\cdot d$,
\[
\det G=s\big[(\alpha\cdot\alpha)(\beta\cdot\beta)-(\alpha\cdot\beta)^2\big]-(d\cdot\alpha)\big[(d\cdot\alpha)(\beta\cdot\beta)-(\alpha\cdot\beta)(d\cdot\beta)\big]+(d\cdot\beta)\big[(d\cdot\alpha)(\alpha\cdot\beta)-(\alpha\cdot\alpha)(d\cdot\beta)\big].
\]
The numerator of $\det Q$ is $N_{kk}N_{vv}-N_{kv}^2$ with $N_{ab}=s(\partial_ad\cdot\partial_bd)-(d\cdot\partial_ad)(d\cdot\partial_bd)$; expanding it term by term gives exactly $s\det G$. Hence $16s^4\det Q=s[d\cdot(\alpha\times\beta)]^2$. \qed

\paragraph{A.17 Curvature bounded by the metric (§\ref{sec:parity}; IQGT).}
For an orthogonal projector $P$ and Hermitian tangents, $Q(V,W)=\Tr(PVW)=\Tr((VP)^\dagger(WP))$, a Frobenius inner product of $VP$ and $WP$. Its Gram matrix on $(V,W)$ is positive semidefinite, so $|Q(V,W)|^2\le Q(V,V)Q(W,W)$, with $Q(V,V)=g(V,V)$ real. Writing $Q(V,W)=g(V,W)+\tfrac i2\Omega(V,W)$, $g(V,W)^2+\tfrac14\Omega(V,W)^2\le g(V,V)g(W,W)$, which is the bound. If $g(V,V)=0$ then $\Omega(V,W)=0$. \qed\ A biorthogonal continuation of $Q$ is a different form and need not keep this positivity (Theorem~\ref{thm:signature}).

\paragraph{What this appendix does not prove.} The numerical statements graded \grade{V} are measurements, and their evidence is the scripts of Appendix C. The two inputs of A.2, the Saha equation, the $1/Z$ coefficients and the Thomas–Fermi law cited in VOTKP, the IQGT relation $g_\parallel\Delta^2=9$ of the gap-step family, the membrane resistivity of the horizon, and the NUBASE2020 and NIST data are outside inputs. Why the census leaves the disk exactly at $\langle K\rangle=2$ on the coupled chain is open (§\ref{sec:status}).
